%% ma: L12-B01-0029
%% lop: 12
%% chuong: 1
%% bai: 1
%% dang: 12.1.2
%% loai: TN4
%% muc: TH
%% dap_an: "D"
%% nguon: "(NBV)-ĐỀ SỐ 18-MỖI NGÀY 1 ĐỀ THI 2025.pdf (D18, MỖI NGÀY 1 ĐỀ - Đề số 18), Phần I Câu 2"
%% kien_thuc: "Bài 1 (tính đơn điệu của hàm số đọc từ đồ thị), Bài 3 (tiệm cận đứng)"
%% trang_thai: chinh-thuc
%% ghi_chu: "Đáp án gốc D đúng. Trùng ý tưởng với dạng đã duyệt 12.1.2 (đọc tính đơn điệu từ đồ thị), nhập chính thức. Đề gốc không cho công thức; hình kho vẽ hàm minh họa y = (x+1)/(x-1) có đồ thị khớp hình gốc (tiệm cận đứng x = 1, tiệm cận ngang y = 1, cắt trục hoành tại x = -1)."
\begin{ex}
Cho hàm số $y=\dfrac{ax+b}{cx+d}$ với $a,b,c,d\in\mathbb R$ có đồ thị như hình vẽ. Mệnh đề nào đúng?
\begin{center}
\begin{tikzpicture}[scale=.8]
  \begin{scope}
    \clip (-3.3,-3.3) rectangle (3.6,3.3);
    \draw[thick,domain=-3.2:.75,samples=80,smooth] plot (\x,{(\x+1)/(\x-1)});
    \draw[thick,domain=1.25:3.5,samples=80,smooth] plot (\x,{(\x+1)/(\x-1)});
  \end{scope}
  \draw[phu] (1,-3.2)--(1,3.2) (-3.2,1)--(3.5,1);
  \draw[truc] (-3.3,0)--(3.6,0) node[below]{$x$};
  \draw[truc] (0,-3.3)--(0,3.4) node[left]{$y$};
  \path (0,0) node[above left]{\footnotesize$O$} (1,0) node[below right]{\footnotesize$1$};
\end{tikzpicture}
\end{center}
\choice{$y'>0,\forall x\ne1$}{$y'>0,\forall x\in\mathbb R$}{$y'<0,\forall x\in\mathbb R$}{\True $y'<0,\forall x\ne1$}
\loigiai{Đồ thị có tiệm cận đứng $x=1$; trên mỗi khoảng $(-\infty;1)$ và $(1;+\infty)$, đồ thị đi xuống (từ trái sang phải), nên hàm số nghịch biến trên từng khoảng đó. Vậy $y'<0$ với mọi $x\ne1$. Chọn D.}
\end{ex}
